\documentclass[12pt]{article}
\usepackage[T1]{fontenc}
\usepackage[utf8]{inputenc}
\usepackage{lmodern}
\usepackage{textcomp}
\usepackage{enumitem}
\usepackage{geometry}
\geometry{margin=1in}
\begin{document}
\begin{center}
Answer Key - Psychology - Undergraduate Level Assessment 21 \
\small (Answers are ordered exactly as in the question paper.)
\end{center}

\section*{Multiple Choice}
\begin{enumerate}[label=\arabic*.]
\item \textbf{Answer:} A
\\ \textit{Options:} A) Episodic; B) Declarative; C) Semantic; D) Working; E) Echoic
\\ \textit{Note:} Episodic memory is responsible for the storage and retrieval of specific personal experiences and events, particularly those that have occurred recently, making it the correct answer for recalling information acquired within the past few hours to days.
\item \textbf{Answer:} A
\\ \textit{Options:} A) identification; B) alliance; C) interpersonal learning; D) disclosure; E) catharsis
\\ \textit{Note:} Identification is the correct answer because, according to Irvin Yalom, it refers to the process by which group members relate to and emulate one another, paralleling the therapeutic bond formed between a therapist and a client in individual therapy.
\item \textbf{Answer:} A
\\ \textit{Options:} A) unethically unless she obtains a waiver from the plaintiff.; B) ethically as long as she did not have a prior relationship with the plaintiff.; C) unethically unless she has a prior relationship with the defendant.; D) ethically as long as she maintains neutrality throughout the trial.; E) ethically only if she has a prior relationship with the defendant.
\\ \textit{Note:} The correct answer is based on the ethical principle that a psychologist must maintain clear boundaries between roles to avoid conflicts of interest, and without a waiver from the plaintiff, acting as both a fact witness and an expert witness compromises the objectivity and integrity of the testimony provided in court.
\item \textbf{Answer:} D
\\ \textit{Options:} A) having the client confront the actual feared object or situation; B) instructing the client to picture each fearful image while maintaining a relaxed state; C) constructing a hierarchy of feared images; D) helping the client experience the desired state of relaxation through hypnosis; E) instructing the client to express their feelings about each feared image
\\ \textit{Note:} The correct answer is accurate because systematic desensitization primarily involves gradual exposure to anxiety-provoking stimuli while teaching relaxation techniques, rather than utilizing hypnosis to achieve relaxation.
\item \textbf{Answer:} A
\\ \textit{Options:} A) Age and gender; B) Birth order and marital status; C) Sleep patterns and dietary habits; D) Astrological sign and time of day; E) Social media usage and television watching habits
\\ \textit{Note:} Age and gender are significant factors influencing cognitive development and socialization processes, which contribute to variations in problem-solving abilities among individuals.
\end{enumerate}

\section*{True / False}
\begin{enumerate}[label=\arabic*.]
\setcounter{enumi}{5}
\item \textbf{Answer:} True
\\ \textit{Options:} A) True; B) False
\\ \textit{Note:} This statement is true because the presence of multiple bystanders often leads to the diffusion of responsibility, where individuals assume someone else will take action, thereby reducing the likelihood of any one person helping in an emergency.
\item \textbf{Answer:} True
\\ \textit{Options:} A) True; B) False
\\ \textit{Note:} The standard deviation is calculated as approximately 5.92 for the sample measurements 1, 3, 7, 10, and 14, based on the formula for sample standard deviation which accounts for both the mean and the variance of the data set.
\item \textbf{Answer:} True
\\ \textit{Options:} A) True; B) False
\\ \textit{Note:} The Wechsler Intelligence Scale for Children-IV includes subtests that effectively measure nonverbal cognitive abilities, making it suitable for assessing general cognitive abilities in nonverbal adolescents.
\item \textbf{Answer:} True
\\ \textit{Options:} A) True; B) False
\\ \textit{Note:} Individuals with Moderate Mental Retardation generally have cognitive abilities that enable them to learn and master academic skills comparable to those of an average eighth grader, allowing for functional literacy and numeracy in everyday situations.
\item \textbf{Answer:} False
\\ \textit{Options:} A) True; B) False
\\ \textit{Note:} The Premack Principle posits that a high-probability behavior can reinforce a low-probability behavior, meaning the reinforcer is not a specific reward but rather the opportunity to engage in a more preferred activity following the desired behavior.
\end{enumerate}

\section*{Long Form Response}
\begin{enumerate}[label=\arabic*.]
\setcounter{enumi}{10}
\item \textbf{Answer:} Psychoactive drugs are substances that induce changes in mental processes, affecting consciousness, mood, perception, and behavior. They can be categorized based on their functions into several types: depressants, which slow down brain activity (e.g., alcohol and benzodiazepines); stimulants, which increase alertness and energy (e.g., caffeine and cocaine); narcotics, which relieve pain and induce euphoria (e.g., morphine and heroin); and hallucinogens, which alter perception and can cause hallucinations (e.g., LSD and psilocybin). Each category serves distinct psychological effects, illustrating the diverse ways in which these drugs influence human experience. Understanding these classifications is crucial for comprehending their implications for mental health and substance use.
\\ \textit{Note:} correct because it accurately defines psychoactive drugs as substances that alter mental processes and categorizes them into four main types-depressants, stimulants, narcotics, and hallucinogens-each with specific functions and examples that illustrate their effects on consciousness and behavior.
\item \textbf{Answer:} Aging significantly impacts sexual functioning, particularly in women, where physiological changes such as decreased estrogen production can lead to alterations in vaginal lubrication. Contrary to common assumptions, some studies suggest that older women may experience increased vaginal lubrication due to hormonal changes and greater familiarity with their bodies, which can enhance sexual pleasure. These physiological changes are often accompanied by psychological implications, including shifts in self-esteem and sexual identity, as societal perceptions of aging can affect sexual confidence. Overall, understanding the relationship between aging and sexual functioning requires a holistic view that considers both the physiological and psychological dimensions, emphasizing that sexual health can continue to thrive in later life.
\\ \textit{Note:} The answer correctly identifies the complex interplay between aging and sexual functioning by addressing the physiological changes, specifically decreased estrogen affecting vaginal lubrication, while also acknowledging the potential for increased lubrication and enhanced sexual pleasure, alongside the psychological implications of these changes.
\end{enumerate}

\vfill
\noindent\textit{End of Answer Key}
\end{document}